\documentclass[12pt]{article}

\usepackage[utf8]{inputenc}
\usepackage[T1]{fontenc}
\usepackage[spanish,es-noshorthands]{babel}
\usepackage{geometry}
\geometry{left=2.5cm, right=2.5cm, top=2.5cm, bottom=2.5cm}
\usepackage{amsmath}
\usepackage{listings}
\usepackage{tikz}
\usetikzlibrary{arrows.meta, positioning, shapes, fit}
\usepackage{hyperref}
\usepackage{setspace}
\usepackage{parskip}
\usepackage{multirow}

% ── estilo de listings ──
\lstdefinestyle{rascalstyle}{
  language={},
  basicstyle=\ttfamily\small,
  breaklines=true,
  breakautoindent=true,
  frame=single,
  framesep=4pt,
  rulecolor=\color{gray!40},
  backgroundcolor=\color{gray!5},
  tabsize=2,
  columns=fixed,
  keepspaces=true,
  extendedchars=true,
  inputencoding=utf8,
  mathescape=true,
  literate={á}{{\'a}}1 {é}{{\'e}}1 {í}{{\'i}}1 {ó}{{\'o}}1 {ú}{{\'u}}1
           {ñ}{{\~n}}1 {Á}{{\'A}}1 {É}{{\'E}}1 {Í}{{\'I}}1 {Ó}{{\'O}}1
           {Ú}{{\'U}}1 {¡}{{!`}}1 {¿}{{?`}}1 {≡}{{$\equiv$}}1
           {→}{{$\rightarrow$}}1 {⇒}{{$\Rightarrow$}}1
           {–}{{--}}1 {°}{{$^\circ$}}1,
}

\lstdefinestyle{kotlinstyle}{
  language=Java,
  basicstyle=\ttfamily\small,
  breaklines=true,
  breakautoindent=true,
  frame=single,
  framesep=4pt,
  rulecolor=\color{gray!40},
  backgroundcolor=\color{gray!5},
  tabsize=2,
  columns=fixed,
  keepspaces=true,
  extendedchars=true,
  inputencoding=utf8,
  literate={á}{{\'a}}1 {é}{{\'e}}1 {í}{{\'i}}1 {ó}{{\'o}}1 {ú}{{\'u}}1
           {ñ}{{\~n}}1 {Á}{{\'A}}1 {É}{{\'E}}1 {Í}{{\'I}}1 {Ó}{{\'O}}1
           {Ú}{{\'U}}1 {¡}{{!`}}1 {¿}{{?`}}1,
}

\lstnewenvironment{rascal}{\lstset{style=rascalstyle}}{}
\lstnewenvironment{kotlin}{\lstset{style=kotlinstyle}}{}

\newcommand{\code}[1]{\texttt{#1}}
\newcommand{\file}[1]{\texttt{#1}}

\begin{document}

%% ── PORTADA ──
\begin{center}
  {\Large\textbf{ISIS-2111 PLE}}\\[0.4cm]
  {\Large\textbf{Proyecto 4}}\\[0.6cm]
  {\LARGE\textbf{VeriLang en Kotlin}}\\[0.8cm]

  Amelia Serrano Arango -- 202221556\\
  Cristian David Ortiz Sanchez -- 202311404\\[0.4cm]
  Mayo 2026
\end{center}

\vspace{0.6cm}

%% ── RESUMEN ──
\section*{Resumen}
\addcontentsline{toc}{section}{Resumen}

En este proyecto se tomó la implementación de VeriLang del Proyecto 3 (con las correcciones
aplicadas de la retroalimentación de la TA) y se adaptó para ejecutarse desde un entorno
Kotlin en lugar de Java. Para ello se creó el módulo \code{RunnerJson.rsc}, que orquesta
el pipeline de análisis (parsing, construcción del AST, verificación semántica y de tipos,
y evaluación de expresiones) y produce un único objeto JSON. Una aplicación de escritorio
en Kotlin con Compose Desktop invoca a Rascal, extrae el JSON resultante y lo despliega
en una interfaz gráfica que muestra el estado del parser, la lista de módulos, los errores
encontrados y la salida del programa. Se añadió soporte para \code{defstruct} y
\code{defdata} en la gramática, el AST, el parser y el checker. El \code{.zip} de entrega
incluye tanto el código fuente de Rascal como la aplicación Kotlin.

\newpage

%% ── TABLA DE CONTENIDO ──
\tableofcontents
\newpage

%% ── 1. INTRODUCCIÓN ──
\section{Introducción}

VeriLang es un lenguaje de programación educativo diseñado en el marco de la materia
ISIS-2111 PLE para explorar conceptos de definición de lenguajes usando Rascal como
\emph{language workbench}. En los proyectos anteriores se definió la gramática formal, se
implementó el parser, el AST, el intérprete y el checker semántico, y se construyó una
interfaz en Java.

Para el Proyecto 4, el objetivo es migrar la interfaz a Kotlin, aprovechando Rascal para
producir un JSON intermedio que la aplicación Kotlin consume. Específicamente se debe:

\begin{itemize}
  \item Reutilizar la implementación de VeriLang del Proyecto 3 con las correcciones
        aplicadas.
  \item Extraer el AST del lenguaje a un archivo JSON mediante un módulo Rascal
        (\code{RunnerJson.rsc}).
  \item Crear un servicio Kotlin (\code{VeriLangService.kt}) que invoque a Rascal,
        capture la salida JSON y la deserialice.
  \item Construir una interfaz gráfica en Compose Desktop (\code{MainWindow.kt}) que
        muestre el estado del parser, la lista de módulos, los errores y la salida.
  \item Entregar un \code{.zip} con la solución y el documento.
\end{itemize}

Adicionalmente, se incorporaron las estructuras \code{defstruct} y \code{defdata}
a la gramática y al sistema de tipos, completando así el conjunto de constructos
solicitados en iteraciones anteriores.

%% ── 2. CORRECCIONES APLICADAS ──
\section{Correcciones aplicadas del Proyecto 3}

Atendiendo la retroalimentación de la TA (Carla) sobre la Entrega 2, se aplicaron las
siguientes correcciones en el Proyecto 3, las cuales se mantienen y heredan en este
Proyecto 4. El detalle completo se encuentra en el documento
\emph{Correcciones aplicadas a VeriLang.pdf}.

\begin{itemize}
  \item \textbf{Identificadores con mayúsculas:} la regla léxica \code{Id} se amplió de
        \code{[a-z]} a \code{[a-zA-Z]} para permitir identificadores como \code{Set},
        \code{Bool} o \code{Nat}.
  \item \textbf{Jerarquía aritmética completa:} se insertaron los niveles
        \code{AddExpr}, \code{MultExpr} y \code{PowExpr} entre \code{EqExpr} y
        \code{UnaryExpr}, añadiendo los operadores \code{+}, \code{-}, \code{*},
        \code{/}, \code{\%} y \code{**}.
  \item \textbf{Operador de equivalencia lógica:} se añadió \code{$\equiv$} (equivalencia) a
        \code{EqExpr}.
  \item \textbf{Literales en Atom:} se conectaron \code{IntLiteral}, \code{FloatLiteral}
        y \code{CharLiteral} como alternativas de \code{Atom}, y se añadieron
        \code{BoolLiteral} y \code{StringLiteral}.
  \item \textbf{CharLiteral con mayúsculas:} se amplió de \code{[a-z]} a \code{[a-zA-Z]}.
  \item \textbf{Evaluador completo:} se implementó un evaluador que reduce expresiones
        a valores concretos (\code{intVal}, \code{floatVal}, \code{boolVal},
        \code{charVal}, \code{stringVal}, \code{appVal}), usando pattern matching
        sobre reglas definidas por el usuario.
\end{itemize}

%% ── 3. GRAMÁTICA FORMAL ──
\section{Gramática formal}

A continuación se presenta la especificación formal actualizada de VeriLang, que incorpora
\code{defstruct} y \code{defdata}.

\subsection{Alfabeto ($\Sigma$)}

\subsubsection{Palabras reservadas}

\begin{lstlisting}[style=rascalstyle, basicstyle=\ttfamily\small]
defmodule  using  defspace  defoperator  defexpression  defrule  defvar
defstruct  defdata  end  forall  exists  in  defer  neg
or  and  true  false  Int  Float  Bool  Char  String
\end{lstlisting}

\subsubsection{Símbolos y operadores}

\begin{lstlisting}[style=rascalstyle, basicstyle=\ttfamily\small]
( ) [ ] : . , + - * / % ** < > <= >= <> = -> =>
or  and  neg  in  ≡
a..z  A..Z  0..9
\end{lstlisting}

\subsection{No terminales ($N$)}

\small
\[
\begin{array}{l}
N = \{\;Module,\;Definition,\;Using,\;SpaceDef,\;SpaceParent,\;StructDef,\;
StructField,\;\\
\qquad DataDef,\;DataVariant,\;DataVariantSignature,\;OperatorDef,\;VarDef,\;
VarDecl,\;\\
\qquad RuleDef,\;ExpressionDef,\;Type,\;LogicalExpression,\;OrExpr,\;AndExpr,\;
EqExpr,\;\\
\qquad AddExpr,\;MultExpr,\;PowExpr,\;UnaryExpr,\;Atom,\;Application,\;\\
\qquad AttributeList,\;Attribute,\;AttributeValue,\;Letter,\;LowerLetter,\;
UpperLetter,\;\\
\qquad Number,\;Identifier,\;IntLiteral,\;FloatLiteral,\;CharLiteral,\;
BoolLiteral,\;StringLiteral\;\}
\end{array}
\]
\normalsize

\subsection{Símbolo inicial ($S_0$)}

\[ S_0 = Module \]

\subsection{Reglas de producción ($P$)}

Para facilitar la lectura, las reglas se presentan agrupadas por categoría. Las
palabras entre comillas dobles son símbolos terminales.

\subsubsection*{Estructura del módulo}

\begin{lstlisting}[style=rascalstyle, basicstyle=\ttfamily\small, mathescape]
Module  ::=  "defmodule" Identifier { Definition } "end"
Using   ::=  "using" Identifier
\end{lstlisting}

\subsubsection*{Definiciones}

\begin{lstlisting}[style=rascalstyle, basicstyle=\ttfamily\small, mathescape, breaklines=false]
Definition     ::=  Using | SpaceDef | StructDef | DataDef
                 |  OperatorDef | VarDef | RuleDef | ExpressionDef

SpaceDef       ::=  "defspace" Identifier [ "<" Identifier ] "end"
StructDef      ::=  "defstruct" Identifier { StructField } "end"
StructField    ::=  Identifier : Type
DataDef        ::=  "defdata" Identifier { DataVariant } "end"
DataVariant    ::=  Identifier [ DataVariantSignature ]
DataVariantSig ::=  : { Type -> }$^+$
OperatorDef    ::=  "defoperator" Identifier : { Type -> }$^+$ [ AttributeList ] "end"
VarDef         ::=  "defvar" { VarDecl "," }$^+$ "end"
VarDecl        ::=  Identifier : Type
RuleDef        ::=  "defrule" Application -> Application "end"
ExpressionDef  ::=  "defexpression" LogicalExpression [ AttributeList ] "end"
\end{lstlisting}

\subsubsection*{Expresiones}

\begin{lstlisting}[style=rascalstyle, basicstyle=\ttfamily\small, mathescape, breaklines=false]
LogicalExpression  ::=  OrExpr
OrExpr             ::=  AndExpr { "or" AndExpr }
AndExpr            ::=  EqExpr { "and" EqExpr }
EqExpr             ::=  AddExpr { ("=" | "<>" | "<" | ">" | "<=" | ">=" | $\equiv$ | "=>" | "in") AddExpr }
AddExpr            ::=  MultExpr { ("+" | "-") MultExpr }
MultExpr           ::=  PowExpr { ("*" | "/" | "%") PowExpr }
PowExpr            ::=  UnaryExpr { "**" UnaryExpr }
UnaryExpr          ::=  [ "neg" ] Atom | "forall" Id "in" Id "." LogicalExpression
                     |  "exists" Id "in" Id "." LogicalExpression
Atom               ::=  Identifier | IntLiteral | FloatLiteral | CharLiteral
                     |  BoolLiteral | StringLiteral | Application | "(" LogicalExpression ")"
Application        ::=  "(" Identifier { LogicalExpression } ")"
\end{lstlisting}

\subsubsection*{Elementos léxicos}

\begin{lstlisting}[style=rascalstyle, basicstyle=\ttfamily\small, mathescape, breaklines=false]
Type           ::=  "Int" | "Float" | "Bool" | "Char" | "String" | Identifier
AttributeList  ::=  "[" { Attribute } "]"
Attribute      ::=  Identifier [ ":" Identifier ]
Identifier     ::=  Letter { Letter | Number | "-" }
Letter         ::=  LowerLetter | UpperLetter
LowerLetter    ::=  "a" | "b" | ... | "z"
UpperLetter    ::=  "A" | "B" | ... | "Z"
Number         ::=  "0" | "1" | ... | "9"
IntLiteral     ::=  Number { Number }
FloatLiteral   ::=  Number { Number } "." Number { Number }
CharLiteral    ::=  "'" Letter "'"
BoolLiteral    ::=  "true" | "false"
StringLiteral  ::=  '"' { cualquier caracter no '"' } '"'
\end{lstlisting}

%% ── 4. ARQUITECTURA ──
\section{Arquitectura}

La comunicación entre Rascal y Kotlin sigue el pipeline ilustrado en la
Figura~\ref{fig:pipeline}. Rascal se encarga de todo el análisis del lenguaje
(parsing, AST, verificación, evaluación) y produce un único objeto JSON. La
aplicación Kotlin invoca a Rascal como un proceso externo, captura el JSON
por salida estándar y lo despliega en la interfaz gráfica.

\begin{figure}[h!]
\centering
\begin{tikzpicture}[
  node distance=1.2cm and 2cm,
  box/.style={rectangle, draw, rounded corners, minimum width=2.8cm,
              minimum height=0.9cm, align=center, font=\small},
  arrow/.style={-{Stealth[scale=1.2]}, thick},
  label/.style={font=\small\itshape}
]

\node[box, fill=blue!10] (vlfile)       {Archivo \texttt{.vl}};
\node[box, fill=green!10, below=of vlfile] (parse)    {Parser\\\texttt{Syntax.rsc}};
\node[box, fill=green!10, below=of parse] (ast)       {AST\\\texttt{Parser.rsc}};
\node[box, fill=orange!10, below=of ast] (check)     {Checker\\\texttt{Checker.rsc}};
\node[box, fill=orange!10, below=of check] (eval)     {Intérprete\\\texttt{Interpreter.rsc}};
\node[box, fill=red!10, below=of eval] (json)      {JSON\\\texttt{RunnerJson.rsc}};

\node[box, fill=purple!10, right=of json] (ktservice) {VeriLang\\Service.kt};
\node[box, fill=purple!10, above=of ktservice] (ktgui)   {Main\\Window.kt};

\draw[arrow] (vlfile) -- node[label, right]{lectura} (parse);
\draw[arrow] (parse)  -- node[label, right]{CST} (ast);
\draw[arrow] (ast)    -- node[label, right]{AST} (check);
\draw[arrow] (check)  -- node[label, right]{AST validado} (eval);
\draw[arrow] (eval)   -- node[label, right]{resultados} (json);
\draw[arrow] (json)   -- node[label, below]{stdout (JSON)} (ktservice);
\draw[arrow] (ktservice) -- node[label, left]{RunResult} (ktgui);

\node[draw, dashed, rounded corners, fit={(parse) (ast) (check) (eval) (json)},
      inner sep=0.6cm, label={[above, font=\bfseries]Rascal}] {};
\node[draw, dashed, rounded corners, fit={(ktservice) (ktgui)},
      inner sep=0.6cm, label={[above, font=\bfseries]Kotlin / JVM}] {};

\end{tikzpicture}
\caption{Pipeline de comunicación Rascal $\rightarrow$ JSON $\rightarrow$ Kotlin.}
\label{fig:pipeline}
\end{figure}

%% ── 5. COMPONENTES ──
\section{Componentes}

\subsection{Syntax.rsc}

Define la gramática libre de contexto de VeriLang usando la notación de
sintaxis concreta de Rascal. Incluye las reglas para \code{defstruct}
y \code{defdata}:

\begin{rascal}
syntax StructDef = structDef:
  "defstruct" Id name StructField* fields "end";
syntax StructField = structField:
  Id name ":" Type tp;

syntax DataDef = dataDef:
  "defdata" Id name DataVariant* variants "end";
syntax DataVariant = dataVariant:
  Id name DataVariantSignature? sig;
syntax DataVariantSignature = dataVariantSignature:
  ":" {Type "-\>"}+ typeSig;
\end{rascal}

Además se añadieron \code{BoolLiteral} y \code{StringLiteral} como
alternativas léxicas, y se actualizó el conjunto \code{Reserved} con
\code{defstruct}, \code{defdata}, \code{true}, \code{false},
\code{Int}, \code{Float}, \code{Bool}, \code{Char}, \code{String}.

\subsection{AST.rsc}

Define las estructuras de datos del árbol de sintaxis abstracta. Se
añadieron los constructores para las nuevas definiciones:

\begin{rascal}
data Definition
  = ...
  | structDefinition(StructDef structDecl)
  | dataDefinition(DataDef dataDecl)
  ;

data StructDef = structNode(str name, list[StructField] fields);
data StructField = structFieldNode(str name, Type tp);
data DataDef = dataNode(str name, list[DataVariant] variants);
data DataVariant = dataVariantNode(str name, list[Type] args);
\end{rascal}

El tipo \code{Literal} unifica todos los literales:

\begin{rascal}
data Literal
  = intLiteral(int intVal)   | floatLiteral(real floatVal)
  | charLiteral(str charVal) | boolLiteral(bool boolVal)
  | stringLiteral(str stringVal)
  ;
\end{rascal}

\subsection{Parser.rsc}

Transforma el árbol de parse concreto (CST) producido por Rascal en el
AST definido en \code{AST.rsc}. Se añadieron las funciones de
transformación para \code{StructDef}, \code{DataDef}, \code{BoolLiteral}
y \code{StringLiteral}. El parser usa pattern matching sobre la sintaxis
concreta de Rascal para cada producción.

Ejemplo para \code{StructDef}:

\begin{rascal}
private AST::StructDef toStructDef(Tree tree) {
  switch (tree) {
    case (StructDef) `defstruct <Id name> <StructField* fields> end`:
      return AST::structNode(text(name),
               [toStructField(field) | field <- fields]);
    default:
      throw "Unsupported struct definition parse tree: <tree>";
  }
}
\end{rascal}

\subsection{Checker.rsc}

El checker combina dos niveles de verificación:

\begin{enumerate}
  \item \textbf{TypePal:} mediante \code{extend analysis::typepal::TypePal},
        se recolectan definiciones y usos de espacios, operadores y variables,
        y se resuelven las restricciones de tipos automáticamente.
  \item \textbf{Validación manual:} se implementaron funciones específicas
        para detectar errores que TypePal no cubre directamente:
        \begin{itemize}
          \item Tipos duplicados (\code{duplicateErrors})
          \item Tipos desconocidos en campos de \code{defstruct},
                argumentos de \code{defdata}, variables y operadores
                (\code{checkTypeExists})
          \item Operadores con aridad insuficiente (\code{checkOperatorShape})
          \item Coherencia de tipos en reglas (\code{checkRule})
          \item Existencia del espacio padre (\code{checkSpaceParent})
        \end{itemize}
\end{enumerate}

El chequeo de tipos de expresiones (\code{typeOf}) verifica cada
constructor del AST:
\begin{itemize}
  \item \textbf{Literales:} infieren su tipo directamente.
  \item \textbf{Variables:} se resuelven contra el entorno de variables.
  \item \textbf{Aplicaciones:} se verifica aridad y tipos de argumentos
        contra la firma del operador.
  \item \textbf{Operadores aritméticos:} requieren \code{Int} o
        \code{Float} con tipo coincidente; \code{\%} solo acepta \code{Int}.
  \item \textbf{Potencia:} base y exponente deben ser \code{Int}.
  \item \textbf{Comparaciones:} \code{=} y \code{<>} requieren tipos
        iguales; \code{$\equiv$} y \code{=>} requieren \code{Bool};
        \code{<}, \code{>}, \code{<=}, \code{>=} requieren \code{Int}
        o \code{Float}; \code{in} verifica compatibilidad con el dominio.
\end{itemize}

\subsection{Interpreter.rsc}

Evalúa las expresiones del AST a valores concretos (\code{RuntimeValue::Val}).
Soporta:

\begin{itemize}
  \item Literales, variables y aplicaciones de operadores.
  \item Reglas definidas por el usuario (\code{defrule}) mediante
        pattern matching (\code{RuleMatcher.rsc}).
  \item Operaciones aritméticas, comparaciones, potencia, negación,
        disyunción y conjunción.
  \item Cuantificadores (\code{forall}, \code{exists}) con entorno
        de variables acotado.
\end{itemize}

La evaluación se realiza sólo si el checker no reporta errores, lo que
garantiza que el programa está bien tipado antes de ejecutarse.

\subsection{RunnerJson.rsc}

Es el módulo central de orquestación. Su función \code{main} recibe la
ruta de un archivo \code{.vl} y ejecuta el pipeline completo:

\begin{enumerate}
  \item Lee el archivo fuente.
  \item Parsing con \code{parse(\#start[Module], src)}.
  \item Construcción del AST con \code{Parser::toAST}.
  \item Verificación con \code{Checker::check} (TypePal + manual).
  \item Evaluación de expresiones con \code{Interpreter::eval} (solo si
        no hay errores).
  \item Impresión de un único objeto JSON con todos los resultados.
\end{enumerate}

Estructura del JSON producido:

\begin{rascal}
{
  "success": true,            // true si todo ok
  "module": "demo",           // nombre del módulo
  "modules": ["demo"],        // lista de módulos
  "parseOk": true,            // estado del parser
  "typeCheckOk": true,        // estado del type checker
  "semanticOk": true,         // estado semántico
  "typeErrors": [],           // errores de tipo
  "semanticErrors": [],       // errores semánticos
  "output": ["resultado"],    // salida de la evaluación
  "error": "",                // mensaje de error general
  "codigoFormateado": "...",  // código fuente original
  "resumen": "modulo=demo; definiciones=9; expresiones=1"
}
\end{rascal}

\subsection{Aplicación Kotlin}

La aplicación de escritorio está construida con Compose Desktop y consta
de cuatro archivos:

\subsubsection*{Main.kt}

Punto de entrada. Crea una ventana con título \texttt{VeriLang} de
800$\times$600 píxeles e inicia \code{MainWindow}.

\subsubsection*{RunResult.kt}

Modelo de datos serializable con \code{kotlinx.serialization}. Sus campos
coinciden exactamente con las claves del JSON producido por
\code{RunnerJson.rsc}:

\begin{kotlin}
@Serializable
data class RunResult(
    val success: Boolean = false,
    val module: String = "",
    val modules: List<String> = emptyList(),
    val parseOk: Boolean = false,
    val typeCheckOk: Boolean = false,
    val semanticOk: Boolean = false,
    val typeErrors: List<String> = emptyList(),
    val semanticErrors: List<String> = emptyList(),
    val output: List<String> = emptyList(),
    val error: String = "",
    val codigoFormateado: String = "",
    val resumen: String = ""
)
\end{kotlin}

\subsubsection*{VeriLangService.kt}

Servicio que ejecuta Rascal como un proceso externo:

\begin{kotlin}
val cmd = listOf(
    "java",
    "-Dfile.encoding=UTF-8",
    "-Drascal.projectPath=${projectRoot.absolutePath}",
    "-jar", rascalJar.absolutePath,
    "RunnerJson",
    filePath
)
val process = ProcessBuilder(cmd).directory(srcDir).start()
\end{kotlin}

Captura el stdout, extrae el primer objeto JSON válido que contenga la
clave \texttt{success}, y lo deserializa en \code{RunResult}. Se ejecuta
en \code{Dispatchers.IO} para no bloquear la interfaz.

\subsubsection*{MainWindow.kt}

Interfaz gráfica que contiene:
\begin{itemize}
  \item Campo de texto para la ruta del archivo y botón \textit{Buscar}
        (usando \code{JFileChooser} filtrado a \code{.vl}).
  \item Botón \textit{Correr} que invoca \code{VeriLangService.run}.
  \item Panel de resultados con indicadores de estado coloreados
        (Parse, Types, Semántica), lista de módulos, resumen del AST,
        errores de tipos, errores semánticos, código formateado y
        salida del programa.
\end{itemize}

\begin{figure}[h!]
\centering
\begin{tikzpicture}[
  box/.style={rectangle, draw, rounded corners, minimum width=6cm,
              minimum height=0.6cm, align=left, font=\small},
  title/.style={font=\bfseries, align=center}
]
\node[title] at (0,1.5) {Distribución de la interfaz};
\node[box, fill=blue!5] (path) at (0,0)     {Ruta: \texttt{/home/.../test.vl}~~[Buscar]~~[Correr]};
\node[box, fill=green!5] (status) at (0,-0.8)
      {[Parse: OK]~~[Types: OK]~~[Sem\'antica: OK]~~m\'odulo: demo};
\node[box, fill=purple!5] (mods) at (0,-1.8) {M\'odulos: demo};
\node[box, fill=yellow!5] (errors) at (0,-2.8) {Errores de tipos: \emph{vac\'io}};
\node[box, fill=orange!5] (output) at (0,-3.8) {Output: \texttt{true}};
\end{tikzpicture}
\caption{Esquema de la interfaz gráfica de MainWindow.kt.}
\label{fig:gui}
\end{figure}

%% ── 6. PRUEBAS ──
\section{Pruebas}

Se incluyen dos archivos de prueba en el directorio \file{instance/}.

\subsection{test.vl (archivo válido)}

\begin{rascal}
defmodule demo
using math
defspace expr end
defspace nat < expr end
defstruct pair
  first : nat
  second : nat
end
defdata truth
  yes
  no
end
defoperator plus : nat -> nat -> nat [commutative associative] end
defoperator leq : nat -> Bool [precedence:low] end
defvar x : nat, y : nat, z : nat end
defrule (plus x y) -> (plus y x) end
defexpression forall x in nat . exists y in nat . neg (x = y) end
end
\end{rascal}

\textbf{Resultado esperado:}
\begin{itemize}
  \item \code{parseOk = true}
  \item \code{typeCheckOk = true}
  \item \code{semanticOk = true}
  \item \code{output} contiene el resultado de evaluar la expresión
        (un valor booleano).
  \item \code{modules = ["demo"]}
\end{itemize}

\subsection{error\_test.vl (archivo con errores)}

\begin{rascal}
defmodule errortest
defspace nat end
defstruct bad-struct
  field : missing-struct-type
  field : nat        // duplicado
end
defdata bad-data
  bad-constructor : missing-constructor-type
  bad-constructor   // duplicado
end
defoperator plus : nat -> nat -> nat end
defoperator leq : nat -> nat -> Bool end
defvar x : nat end
defrule (plus x y) -> (minus x y) end  // minus no existe
defexpression neg (x = z) end           // z no definida
defexpression (undefined-op x y) end    // operador no existe
defvar bad : missing-space end          // tipo no existe
defspace child < missing-parent end     // padre no existe
defexpression forall n in missing-domain . (leq n x) end
defexpression (plus x true) end         // tipo incorrecto
end
\end{rascal}

\textbf{Resultado esperado:}
\begin{itemize}
  \item \code{parseOk = true}
  \item \code{typeCheckOk = false}
  \item \code{semanticOk = false}
  \item \code{typeErrors} contiene al menos 8 errores (campos duplicados,
        tipos desconocidos, constructores duplicados, operador no definido,
        variable no definida, espacio faltante, padre faltante, dominio
        faltante, tipo incorrecto en argumento).
  \item \code{output = []} (no se evalúa si hay errores).
\end{itemize}

%% ── 7. EJECUCIÓN ──
\section{Ejecución}

\subsection{Requisitos}

\begin{itemize}
  \item \textbf{Rascal:} \code{rascal-shell-stable.jar} ubicado en el
        directorio \code{verilang/} o en el directorio padre del proyecto.
  \item \textbf{Kotlin:} Gradle instalado (\code{gradle --version}).
\end{itemize}

\subsection{Comandos}

Desde la terminal, ubicarse en \code{verilang/kotlin-app} y ejecutar:

\begin{kotlin}
gradle run
\end{kotlin}

La aplicación invoca internamente:

\begin{kotlin}
java -Dfile.encoding=UTF-8 \\
     -Drascal.projectPath=<verilang> \\
     -jar rascal-shell-stable.jar \\
     RunnerJson <archivo.vl>
\end{kotlin}

La primera ejecución descarga las dependencias de Compose
($\sim$2 minutos). Las siguientes son inmediatas.

\subsection{Ejemplo}

\begin{enumerate}
  \item Abrir la aplicación con \code{gradle run}.
  \item Presionar \textit{Buscar} y seleccionar \code{instance/test.vl}.
  \item Presionar \textit{Correr}.
  \item Los indicadores \textit{Parse}, \textit{Types} y \textit{Semántica}
        se muestran en verde (OK).
  \item El módulo \code{demo} aparece en la lista de módulos.
  \item La salida de la evaluación se muestra en la sección \textit{Output}.
\end{enumerate}

%% ── 8. CONCLUSIONES ──
\section{Conclusiones}

En este proyecto se logró migrar exitosamente la interfaz de VeriLang de
Java a Kotlin, demostrando la flexibilidad de Rascal como \emph{language
workbench} para separar el análisis del lenguaje de la presentación de
resultados.

\subsection*{Logros principales}
\begin{itemize}
  \item El pipeline Rascal $\rightarrow$ JSON $\rightarrow$ Kotlin funciona
        de extremo a extremo, permitiendo que cualquier lenguaje definido
        en Rascal se conecte a una GUI Kotlin con solo implementar
        \code{RunnerJson.rsc}.
  \item La gramática de VeriLang se completó con \code{defstruct} y
        \code{defdata}, cubriendo todos los constructores solicitados
        durante el semestre.
  \item El checker combina TypePal para resolución automática de tipos
        con validación manual para reglas específicas del lenguaje,
        resultando en un sistema de tipos robusto.
\end{itemize}

\subsection*{Limitaciones}
\begin{itemize}
  \item Los errores de tipo y semánticos se reportan en una sola lista,
        sin diferenciación. Una mejora futura sería separarlos para
        dar retroalimentación más precisa al usuario.
  \item El intérprete maneja cuantificadores de forma simplificada
        (asigna \code{appVal} a las variables cuantificadas), lo que
        limita la expresividad de la evaluación.
  \item La aplicación requiere Gradle instalado; no se incluyen los
        scripts \code{gradlew} por restricciones de tamaño del correo
        corporativo.
\end{itemize}

\subsection*{Trabajo futuro}
\begin{itemize}
  \item Separar los errores del checker en categorías (tipo vs. semánticos)
        para mejorar la claridad en la interfaz.
  \item Implementar un pretty printer que produzca código formateado
        a partir del AST.
  \item Agregar resaltado de sintaxis (\code{@category}) en la gramática
        de Rascal, que no pudo incluirse por conflictos con el sistema
        de palabras reservadas.
  \item Extender el intérprete para manejar semántica más compleja
        (evaluación de \code{defstruct}, \code{defdata} y constructores
        con argumentos).
\end{itemize}

\end{document}
